\section{Discussion}\label{sec:discussion}

\paragraph{One energy, three roles}
In machine learning, an energy-based model scores the compatibility of an input $X$ and an
output $Y$ with a scalar energy, lowest for compatible pairs, and is used, among other
things, to find the $Y$ most compatible with $X$, to rank two candidates $Y_1$ and $Y_2$,
and to detect whether a given $Y$ is compatible with $X$ at all \cite{LeCun2007EBM}. Its
energy is learned and measured in arbitrary units, and detection compares it with a
threshold that is generally unknown when the model is built. For linear elastostatics the
discrete potential energy is such an energy for the problem $(K,F)$, but assembled rather
than learned, and, with compatibility measured in the stiffness norm, it answers the three
questions exactly. Its minimiser is the finite-element solution, which the label-free
surrogate learns to approximate for a whole family of problems by minimising the energy
itself (\cref{cor:training}): the training objective. Differences of energy rank the
predictions of one load case by their stiffness-norm errors, which by \cref{prop:stress}
are compliance-weighted stress errors (\cref{cor:rank}): the evaluation metric. The sign of
the energy answers the third question if a compatible prediction is one no further from the
solution than the zero field; the threshold is then the energy of that field, zero, and
needs no tuning (\cref{cor:zero}): the failure detector, which the rescaling by $\cstar$
completes with a repair that never increases the stiffness-norm error (\cref{prop:amp}).
The check and the repair need only the assembled $K$ and $F$ and apply to every surrogate,
however it was trained. In three dimensions they flagged, and repaired, every load case on
which the supervised transformer or the graph network was worse than the zero field, as
\cref{cor:zero,prop:amp} guarantee (\cref{sec:results3d:detection}).

\paragraph{The norm of the loss}
By \cref{cor:training,prop:stress}, the label-free objective is, up to a constant, the
compliance-weighted mean-square error of the stress field. The supervised transformer and
the graph network were trained on the relative Euclidean error of the displacement, which
weights all eigenmodes of the stiffness matrix equally and so under-weights, relative to the
stiffness norm, the stiff modes that carry the strains (\cref{sec:theory:metrics}). A
supervised network can therefore reach a small displacement error with an error rich in
stiff modes, and a large stress error. Both comparisons tested this explanation by changing
the norm of the supervised loss and nothing else. In two dimensions the von Mises error fell
by more than half, on the validation and on the held-out instances, and the errors, of about
the same size, lay lower in the spectrum
(\cref{sec:results2d:norm,sec:results2d:spectrum,sec:holdout}). The three-dimensional run
repeated the test at the scale of \cref{sec:results3d}: the stiffness-norm loss lowered the
supervised transformer's median von Mises error by \numKdVmRel{} and its median relative
energy gap by \numKdGapRel{}, on the seeds' medians, a statistic chosen after the
per-instance values of the networks compared were known, and on these medians it removed
\numKdClosureVm{} and \numKdClosureGap{} of the supervised transformer's excess over the
label-free transformer in the two errors, a descriptive reading (\cref{sec:results3d:norm}).
The label-free transformer's advantage over the supervised one thus lies in the norm of its
loss rather than in the absence of labels. With labels, training in the stiffness norm may
do better still: in three dimensions its median relative energy gap was \numKdFreeRel{}
below the label-free transformer's, beyond the guard, though with three seeds the Welch 95\%
interval runs from \numKdFreeWelchLow{} to \numKdFreeWelchHigh{}; in two dimensions no
difference beyond the guard was shown, on the validation instances or on held-out ones; on
the validation instances the stiffness-norm transformer was the more accurate in every seed,
but on the held-out ones its seed-mean relative energy gap was the higher
(\cref{sec:results2d:norm,sec:holdout}). That comparison changes the weighting of the load
cases \eqref{eq:lkgrad} and the learning rate together with the labels, and which of them
accounts for the difference was not tested. The label-free objective is thus not necessarily
the most accurate way to train in the stiffness norm when labels are at hand; it is the way
to train in it without them. In three dimensions the supervised transformer's errors lay
higher in the spectrum than the label-free transformer's on most load cases
(\cref{sec:results3d:reverse}). The energy term of the supervised transformer is a partial
remedy: it added the stress-weighted signal and removed most of the large failures, but on
the same instances the label-free transformer, which receives only that signal, was more
accurate still (\cref{sec:results3d:same}). Few-shot training on displacement labels behaved
in the same way on the finer mesh, restoring the displacement error but not the energy and
stress errors (\cref{sec:transfer:zeroshot}). The test of the mechanism rests on one model
size in each dimension and on readings taken after the runs; it supports the claim of
\cref{sec:intro} rather than establishing it in general.

\paragraph{A classical baseline}
Linear elastostatics is a test bed in which the identity can be proved, not a regime in
which networks are needed. With the stiffness matrix assembled, a polynomial Ritz projection
of a few hundred unknowns, a baseline added after the original analyses, took a fraction of
a second per validation instance; at its highest degree its displacement error was smaller
than each three-dimensional network's, in each seed, on at least \numRitzValNineDispAll{} of
the instances, and its relative energy gap was comparable to the label-free transformer's
(\cref{sec:results3d:ritz}). The networks' advantages lie elsewhere: the projection's median
peak stress error was \numRitzValPeakRatioMin{} to \numRitzValPeakRatioMax{} times the
label-free transformer's and larger than the stiffness-norm transformer's on every instance,
and no network needs a stiffness matrix at inference. The identity does not depend on the
approximation: the energy gap evaluates a Ritz solution, and the sign of the energy screens
it, as they do a network's prediction. Neural predictions have also been used to enrich
finite-element approximation spaces \cite{Barucq2026Enrich}. The rescaling by $\cstar$ is
the smallest such enrichment, the projection onto the span of one prediction
(\cref{prop:amp}); combining a prediction with a polynomial space we leave to separate work.

\paragraph{A counterexample}
One comparison, reported in a preprint, points the other way. In the pressure vessels of
Sun et~al.~\cite{Sun2026FEVessel}, the same point-cloud transformer trained on the relative
error of the displacement labels instead of the energy, with every other setting unchanged,
was more accurate in displacement and in the $H^1$ seminorm (their Table 2), against the
three-dimensional result of \cref{sec:results3d} and against the account above. The two
studies differ in ways that
may matter. That comparison is a single run, and its reported errors are the lowest test
errors over the last third of training. Its reference solutions use reduced-integration
hexahedra with hourglass control, and its energy loss takes derivatives from the shape
functions of those elements, with a quadrature that is not stated. \Cref{lem:exact} makes
energy training exact regression on the reference solution only when the energy minimised is
the assembled energy of the discretisation that defines the reference; if the two differ,
the minimiser of the energy loss is not the reference solution, and energy training
regresses towards a different target. Whether this explains the reversal is not known. The
comparison shows that the advantage of label-free training reported here is not universal.

\paragraph{The graph network}
In both dimensions the graph network had the lowest displacement error, in two dimensions
within the noise guard of the supervised transformer's, and von Mises errors several times
the label-free transformer's (\cref{sec:results3d:reverse,sec:results2d:accuracy}). Its
errors lay highest in the spectrum of $K$ (\cref{sec:results3d:reverse,sec:results2d:spectrum}):
relatively more weight in stiff modes, as for an error that varies rapidly from node to
node, which \cref{fig:field2d,fig:field-worst,fig:field-median} show as rough stress fields.
We did not investigate why message passing over the mesh edges produces such errors. The
observation matters for evaluation more than for the graph network itself: a ranking of
surrogates by displacement error alone can invert their ranking by stress, the quantity on
which design decisions usually depend.

\paragraph{Warm starts}
The label-free prediction did not shorten conjugate gradients to a relative residual of
\numKfiveTol{} in two dimensions (\cref{sec:cost:warm}), nor to \numCostCgTol{} in two or
three dimensions (\cref{tab:cost}), whereas the pretrain finite element
method needed up to about nine times fewer conjugate-gradient iterations than a zero start
at a tolerance of $10^{-3}$ \cite{Wang2026PretrainFEM}. In its pressure-vessel follow-up the
saving of conjugate gradients fell from about three times at a relative residual of
$10^{-3}$ to almost none at $10^{-4}$ \cite{Sun2026FEVessel}, and neural-operator warm
starts have also been reported to accelerate iterative solvers \cite{Eshaghi2026NOWS}.
Wang et~al.~\cite{Wang2026PretrainFEM} and Sun et~al.~\cite{Sun2026FEVessel} explain why
the benefit shrinks as the tolerance tightens: a warm start lowers an iteration bound by a
fixed number of iterations, as Wang et~al.\ show for fixed-point iterations and Sun et~al.,
in the relative residual, also for conjugate gradients. The bound is not the count: with
conjugate gradients, the warm start of Sun et~al.\ saved 580 iterations at $10^{-3}$ but 69
at $10^{-4}$ (their Table 3). \Cref{cor:warm} states the bound for conjugate gradients in the
energy norm: a warm start lowers the iteration bound by a number of iterations set by its
own accuracy, so the share of the run it can save shrinks as the tolerance tightens, and
a residual-based stop can hide part of it (\cref{rem:residual}). Reported warm-start
speed-ups are therefore comparable only when the tolerance, the stopping test and the
preconditioner are stated with them: Wang et~al.~\cite{Wang2026PretrainFEM} report, for
their linear problems, iteration counts to a relative residual of $10^{-3}$, with
unpreconditioned conjugate gradients in two dimensions and with unpreconditioned and
preconditioned conjugate gradients in three, and Sun et~al.~\cite{Sun2026FEVessel} report
iteration counts and wall-clock speed-ups at three relative-residual tolerances, the
speed-ups computed from times that include the network's forward pass and, for algebraic
multigrid, its set-up.

\paragraph{Transfer}
No network transferred to the finer mesh, and the readings of \cref{sec:transfer} point at
the output scaling: the network's output is multiplied by the largest nodal load, which
depends on the mesh, while the solution does not. The networks of the July run, whose output
was not scaled in this way, lost little accuracy on meshes \numCoarsenLow{} and
\numCoarsenHigh{} times finer than those they were trained on, although that evidence is
weak (\cref{sec:transfer:amplitude}). Whether an output scale independent of the mesh
removes the failure is the subject of separate work; until it is shown to, the label-free
transformer of this paper should be regarded as valid within the range of mesh sizes it was
trained on. The Ritz projection, which has no such scale, kept most of its accuracy on the
finer mesh, though not in peak stress (\cref{sec:transfer:zeroshot}).

\paragraph{Limitations}
The study is restricted to linear elastostatics, where \cref{lem:exact} is exact; one family
of geometries and one set of four load cases per dimension; and three seeds per network.
\Cref{prop:stress}, which ties the stiffness norm to the stresses, assumes a homogeneous
isotropic material and an exactly integrated stiffness, as with the linear elements used
here; and the reported von Mises error weights elements equally, so \cref{prop:stress}
explains it rather than bounds it, although the area- and volume-weighted errors, measured
after the runs, differed little (\ref{app:repro:checks}). The displacement comparison
reversed between the dimensions, with the same code
(\cref{sec:results3d:same,sec:results2d:accuracy}), and we have not identified the cause.
The two-dimensional comparison was evaluated on the 256 validation instances on which the
results that motivated it had been seen, and the values of its label-free networks, and the
directions of its hypotheses at a far smaller scale, were known when it was registered
(\ref{app:prereg}); its comparisons across numbers of labelled instances are exploratory.
Its claims, and those of the three-dimensional comparison, were read again on held-out
instances, where seven of eight replicated and the reduction of the relative energy gap by
the stiffness-norm loss in two dimensions did not (\cref{sec:holdout}). The stiffness-norm
control was run once in each dimension, at one model size and with the supervised settings,
not tuned for it; in three dimensions it was trained about seven weeks after the network it
is compared with, and its verdicts read the seeds' medians, a statistic chosen after the
per-instance values of the networks it compares were known (\cref{sec:results3d:norm}). The
spectral readings that bear on its mechanism were taken after the two-dimensional run, and
the Ritz baseline was added after the original analyses, with its specification written and
made public before its paired computations (\cref{sec:methods:ritz}). The scope of this
paper, the energy objective without the joint-embedding term that a pre-registered
comparison kept, was chosen after that comparison, and the three-dimensional gate was
amended after the initial run; \ref{app:prereg} reports both. The pre-registration documents
of the two runs of July 2026 are weaker records than the later ones
(\cref{sec:methods:prereg}). No network transferred to a substantially finer mesh; the
discretisation invariance reported for point-cloud operators trained on the energy
\cite{Wang2026PretrainFEM,Sun2026FEVessel} was tested by varying the density of the input
point cloud rather than by refining the mesh, and label-free fine-tuning on the finer mesh,
which needs no solves, was not tested here.

\paragraph{Outlook}
The construction extends wherever each step of a computation is a convex minimisation. In
linear dynamics, every implicit step of a Newmark-type integrator solves a symmetric
positive definite system whose effective stiffness combines the stiffness, mass and
damping matrices, and \cref{lem:exact} holds step by step \cite{CaoSong2026}. Incremental
variational formulations of dissipative materials write each step of a nonlinear history as
the minimisation of an incremental potential \cite{OrtizStainier1999}, and variational
integrators derive each step of a time integration from a discrete variational principle
\cite{MarsdenWest2001}. Where such a step is a minimisation, its potential remains a
label-free training objective, although the energy gap is no longer a squared norm, and the
counterparts of \cref{prop:amp,cor:zero} remain to be worked out; at finite strains the
incremental potential is in general not convex \cite{OrtizStainier1999}, so that its
minimiser need not be unique. These are the problems in which a solve may be expensive
enough for label-free training to pay for itself.
